\documentclass[11pt]{article}
\usepackage[margin=1in]{geometry}
\usepackage{amsmath,amssymb,amsthm,hyperref}
\setlength{\emergencystretch}{2em}
\newtheorem{theorem}{Theorem}
\title{A singleton spectral boundary is not a semicircle\\Disproof of Conjecture 00000007157}
\author{ziangni-sys}
\date{October 8, 2026}
\begin{document}
\maketitle

\paragraph{Statement and interpretation.}
Both supplied language versions assert that the spectral boundary of a
nonnormal complex matrix is a semicircle whose radius is its numerical
radius. We use ``spectral boundary'' in its standard literal sense, the
topological boundary of the spectrum in $\mathbb C$. The counterexample
below rules out every center and orientation of the proposed semicircle.
It does not refute an alternative statement giving a semicircular enclosure
or an inequality: such a statement is different from saying that the
boundary \emph{is} a semicircle. No enclosure formula or additional
hypothesis is specified in the source.

\begin{theorem}
The nonnormal matrix
\[
                J=\begin{pmatrix}0&1\\0&0\end{pmatrix}
\]
has spectral boundary $\partial\sigma(J)=\{0\}$ and strictly positive
numerical radius. Consequently its spectral boundary is not a semicircle
of its numerical radius.
\end{theorem}
\begin{proof}
First,
\[
 JJ^*=\begin{pmatrix}1&0\\0&0\end{pmatrix},\qquad
 J^*J=\begin{pmatrix}0&0\\0&1\end{pmatrix},
\]
so $J$ is nonnormal. For any complex $z$,
\[
          \det(zI-J)=\det\begin{pmatrix}z&-1\\0&z\end{pmatrix}=z^2.
\]
Thus $zI-J$ is invertible exactly when $z\ne0$, and $\sigma(J)=\{0\}$.
A singleton in $\mathbb C$ is closed and has empty interior, hence its
topological boundary is itself.

For a unit vector $v=(v_1,v_2)\in\mathbb C^2$ with the Euclidean norm,
\[
      v^*Jv=\overline{v_1}v_2,
      \qquad |v_1|^2+|v_2|^2=1.
\]
Therefore the numerical radius, by its actual quadratic-form definition, is
\[
 w(J)=\sup_{|v_1|^2+|v_2|^2=1}|\overline{v_1}v_2|.
\]
The inequality $(|v_1|-|v_2|)^2\geq0$ gives
$|\overline{v_1}v_2|=|v_1||v_2|\leq1/2$. Conversely, the rational vector
$v=(3/5,4/5)$ is a unit vector and gives
$|\overline{v_1}v_2|=12/25$. Hence
\[
                  0<\frac{12}{25}\leq w(J)\leq\frac12.
\]
These bounds, rather than an assumed numerical-radius formula, suffice.

Let $r=w(J)>0$. Any closed semicircle of radius $r$ can be expressed as
\[
 S(c,u,r)=\{c+u\zeta:\ |\zeta|=r,\ \operatorname{Im}\zeta\geq0\},
       \qquad c\in\mathbb C,\quad |u|=1.
\]
Both $c+ur$ and $c-ur$ belong to this semicircle, and they are distinct
because $u\ne0$ and $r>0$. They cannot both belong to $\{0\}$.
Thus $\partial\sigma(J)\ne S(c,u,w(J))$ for every center and orientation.
The same elementary cardinality obstruction also applies if ``semicircle''
excludes its endpoints: a positive-radius open semicircular arc still
contains more than one point.
\end{proof}

\paragraph{Formal verification.}
The Lean project represents $J$ as a genuine $2\times2$ complex matrix,
and checks nonnormality using conjugate transpose and matrix multiplication.
Its spectrum is Mathlib's algebraic spectrum: invertibility of $zI-J$ is
reduced to its determinant, establishing \texttt{spectrum\_singleton}.
The standard topological frontier yields \texttt{boundary\_singleton}.
The formal numerical radius is the real supremum of the moduli of
$v^*Jv$ over all Euclidean unit vectors, with the Euclidean unit condition
written explicitly in coordinates. The proof establishes nonemptiness,
boundedness, and the bounds $12/25\leq w(J)\leq1/2$.
The final theorem \texttt{boundary\_not\_semicircle} quantifies over arbitrary
complex centers and unit complex orientations. It uses the actual
positive numerical radius, not a scalar substitute or an assumed
spectral certificate. This directly negates the conjecture's universal
boundary assertion.
\end{document}
